% ---------------------------------------------------------------
% Terminologia específica
% ---------------------------------------------------------------
\chapter*{Terminologia específica}
\addcontentsline{toc}{chapter}{Terminologia específica}

Termos técnicos usados neste relatório, por ordem alfabética.

\begingroup
\small
\rowcolors{2}{white}{corFundo}
\begin{longtable}{@{}L{0.25\textwidth}L{0.70\textwidth}@{}}
\linhacab \cab{Termo} & \cab{Definição} \\
\endfirsthead
\linhacab \cab{Termo} & \cab{Definição} \\
\endhead
Arbitragem & Escolha, no fim de cada ciclo, da decisão de velocidade e da decisão de via entre todas as propostas das regras. \\
Base de conhecimento & Conjunto das regras que representam o conhecimento do domínio, separado do motor que as aplica. \\
Encadeamento progressivo & Inferência que parte dos factos e aplica as regras cujas condições se verificam (\ing{forward chaining}). \\
Encadeamento regressivo & Inferência que parte de um objetivo e procura os factos que o provam (\ing{backward chaining}). \\
\ing{Entry point} & Em \ing{Drools}, fluxo com nome pelo qual os eventos de um tipo entram na sessão. \\
Explicação \ing{Why} / \ing{How} / \ing{Why Not} & Explicações de porque foi tomada uma decisão, de como se chegou a ela e de porque não foi tomada outra. \\
Facto & Unidade de informação sobre a situação atual, mantida na memória de trabalho. \\
Lógica difusa & Raciocínio com categorias vagas, a que um valor pertence com um grau entre 0 e 1 (\ing{fuzzy logic}). \\
Manutenção da verdade & Retirada automática de uma conclusão quando o facto que a justificava deixa de existir (\ing{truth maintenance}). \\
Memória de trabalho & Conjunto dos factos que o sistema conhece num dado momento (\ing{working memory}). \\
Metaconhecimento & Conhecimento sobre as próprias regras, que controla a inferência: ordem das fases, prioridades, desempate. \\
Motor de inferência & Componente que avalia as regras sobre a memória de trabalho e produz conclusões. \\
Nível de risco & Gravidade de uma situação: baixo, moderado, alto ou crítico. Acompanha a decisão, mas não a seleciona. \\
Operadores de Allen & Relações entre instantes e intervalos de tempo, como \ing{after} ou \ing{coincides}, usadas nas regras temporais. \\
Prioridade & Categoria, de 10 a 70, que decide entre propostas concorrentes. \\
Proposta & Ação sugerida por uma regra num ciclo, com velocidade, risco, prioridade, razão e evidência. \\
Pseudo-relógio & Relógio da sessão \ing{Drools} avançado pelo programa, exatamente 100 ms por ciclo. \\
Saliência & Em \ing{Drools}, ordem relativa de disparo das regras ativas (\ing{salience}). \\
Sessão com estado & Sessão do motor que mantém os factos entre ciclos (\ing{stateful KieSession}). \\
Sistema pericial & Sistema que representa conhecimento especializado de forma explícita e o utiliza, através de um motor de inferência, para chegar a conclusões justificáveis (\ing{expert system}). \\
\ing{Snapshot} & Estado completo do sistema num ciclo, enviado à interface. \\
Teto de velocidade & Proposta que diz «não mais depressa do que\ldots»; todos os tetos se aplicam e vence o mais baixo. \\
TTC & Tempo até à colisão (\ing{time to collision}): distância ao veículo da frente a dividir pela velocidade de aproximação. \\
\ing{WebSocket} & Protocolo de comunicação persistente e bidirecional sobre uma única ligação TCP (RFC 6455). \\
\end{longtable}
\addtocounter{table}{-1} % a tabela da terminologia não tem número (como no Word)
\endgroup
\clearpage
